% S4: chi duoc \input boi paper2_hujos/supplementary.tex (kiem tra 2026-10-09).
% Truoc day la \section{...} danh so tu dong -> render trong supplementary thanh
% "1 Subsidence Invariance" trong khi than bai trich dan "S4" (peer review F5).
\section*{S4. Subsidence Invariance}

AWD compliance is usually stated as an absolute water-table or ponding depth,
for example $15$\,cm below the surface, the reflooding threshold of the Safe-AWD
treatment in~\cite{balaine2019}. In a subsiding delta
the surface elevation $z_s(t)$ falls at rate $\dot z$. Two measurement protocols
behave differently and must be kept apart: if the threshold is referenced to a
\emph{fixed elevation datum} then datum and surface move apart and the effective
threshold drifts, whereas if it is measured \emph{relative to the current field
surface}, as with a perched-water-table tube set from the surface, datum and surface
fall together and the threshold does not drift. The proposition below concerns the
first protocol only.

\begin{proposition}[Drift of a datum-referenced threshold]
\label{prop:subsid}
Consider a threshold expressed against a \emph{fixed elevation datum}, such as a
water level read from a sensor anchored to a survey benchmark. Over a crediting
period of $T_c$ years that threshold shifts by $\dot z\,T_c$ relative to the field
surface. For $\dot z\in[1.1,2.5]$\,cm/yr and $T_c=5$\,yr the shift is
$5.5$--$12.5$\,cm, that is $37$--$83\%$ of a $15$\,cm threshold. A criterion
written in the depletion fraction \eqref{eq:depl} has zero drift under either
protocol, because $FC$ and $WP$ are soil properties referenced to the root zone
and not to an absolute elevation.
\end{proposition}

The scope of the proposition should be stated plainly. It does \emph{not} apply to
a threshold measured relative to the current field surface: the $15$\,cm
below-surface criterion of the Safe-AWD treatment in~\cite{balaine2019} is a
perched water table measured from the surface, so it translates with the surface
and does not automatically drift by the full subsidence amount. We could not verify
from that source which tube type the deployed sites use, so we assign no protocol
to them; the $5.5$--$12.5$\,cm shift is an upper bound valid for a datum-anchored
sensor.

\begin{proof}
The datum-referenced criterion compares a measured depth to a constant benchmark;
if the benchmark and the surface move apart at $\dot z$, the effective threshold
changes by $\dot z T_c$ over $T_c$ years. If the benchmark is reset to the current
surface the two translate together and the drift is zero, which is why the
proposition is stated for the first protocol only. The depletion criterion compares
$FC-w_t$ to $(FC-WP)$, both of which are defined within the root zone and
translate rigidly with the surface, so their ratio is invariant to $z_s$.
\end{proof}

The subsidence rates are those measured for the delta by
\cite{minderhoud2017}. The practical consequence is that a scheme anchored on
depletion fraction needs no re-surveying of field benchmarks during a five-year
crediting period, whereas a scheme reading from a datum-anchored sensor must have that datum re-surveyed
periodically; a surface-relative measurement needs neither. In both cases the
depletion fraction is the quantity that requires no elevation correction.
